Fix \(n\) and \(P\in\mathcal P_{\mathrm{NG},n}\).  All norms in the
first part are with respect to the probability law \(P_X\).  The code
errors are measurable and bounded, so all displayed norms are finite,
including when \(s=\infty\).

If \(r\le s<\infty\), Jensen's inequality for the concave function
\(u\mapsto u^{r/s}\) gives, for every measurable \(f\),
\[
 E_P[|f(X)|^r]
 =E_P\!\left[(|f(X)|^s)^{r/s}\right]
 \le\left(E_P[|f(X)|^s]\right)^{r/s}.
\]
Taking \(r\)-th roots yields
\[
 \|f\|_{L^r(P_X)}\le\|f\|_{L^s(P_X)}.                 \tag{1}
\]
For \(s=\infty\), the almost-sure inequality
\(|f(X)|\le\|f\|_{L^\infty(P_X)}\), followed by integration, proves
(1) again.  The assumed \(s\ge r+1\) therefore suffices.

Apply (1) to \(f=\bar g_n-g_0\) and \(f=\bar q_n-q_0\).  The two
radius conditions in \(\text{def:non-gaussian-class}\) give
\[
 \|\bar g_n-g_0\|_{L^r(P_X)}\le\varepsilon_{1,n},
 \qquad
 \|\bar q_n-q_0\|_{L^r(P_X)}\le\varepsilon_{2,n}.       \tag{2}
\]
These are exactly the code restrictions in
\(\text{def:jms-ace-class}\).  Every other member condition in that
definition---the conditional-mean restriction, the target and nuisance
range restrictions, the two sub-Gaussian restrictions, independence of
\(\eta\) and \(X\), and the order-\(r+1\) cumulant separation---also
appears in \(\text{def:non-gaussian-class}\).  Hence
\[
 \mathcal P_{\mathrm{NG},n}
 \subseteq\mathcal P_{\mathrm{ACE},n}^{\mathrm{JMS}}.   \tag{3}
\]
This is not an equality assertion.  For example, on a probability space
with an event \(A\) satisfying \(0<P_X(A)<1\), the bounded function
\(f=c\mathbf 1_A\), \(c\ne0\), obeys
\[
 \frac{\|f\|_{L^s(P_X)}}{\|f\|_{L^r(P_X)}}
 =P_X(A)^{1/s-1/r}>1
\]
when \(s<\infty\); for \(s=\infty\), the ratio is
\(P_X(A)^{-1/r}>1\).  Thus the reverse norm implication at the same
radius does not follow.

It remains to establish the published comparator conclusion.  Assume
\(\mathsf E_n^{\mathrm{JMS}}\).  By
\(\text{def:jms-eligibility}\), the logarithms and denominators in the
eligibility display are defined.  In particular, \(\varepsilon_{1,n}>0\).
If \(\varepsilon_{2,n}>0\), the cited result
\(\text{lem:jms-ace-quantile-upper}\) gives exactly
\[
 \sup_{P\in\mathcal P_{\mathrm{ACE},n}^{\mathrm{JMS}}}
 Q_{P,1-\gamma}
 \!\left(|\widehat\theta_{\mathrm{ACE},r,n}-\theta_0(P)|\right)
 \le B_{\mathrm{ACE},n,\gamma}.                         \tag{4}
\]
If \(\varepsilon_{2,n}=0\), choose
\(u\in(0,\varepsilon_{1,n})\) and enlarge the outcome-code radius at
index \(n\) from zero to \(u\).  The original uncertainty class is a
subset of this enlarged class.  Because \(u<\varepsilon_{1,n}\), the
maximum in \(b_{2,n}\) is unchanged, and hence the same eligibility
condition holds.  The cited positive-radius bound applied to the enlarged
class gives
\[
 \begin{aligned}
 &\sup_{P\in\mathcal P_{\mathrm{ACE},n}^{\mathrm{JMS}}}
 Q_{P,1-\gamma}
 \!\left(|\widehat\theta_{\mathrm{ACE},r,n}-\theta_0(P)|\right)\\
 &\quad\le C_{\gamma}r!16^r\delta^{-1}
 \Big[\varepsilon_{1,n}^r u+C_{\theta}\varepsilon_{1,n}^{r+1}
 +64(C_g+\psi_{\eta})^r
 \{r^2(C_g+\psi_{\eta})+\psi_{\xi}+C_{\theta}\psi_{\eta}\}
 (\gamma n)^{-1/2}\Big].
 \end{aligned}
\]
Letting \(u\downarrow0\) proves (4) also at the zero-radius boundary.
By definition, the right-hand side of (4) is
\(B_{\mathrm{ACE},n,\gamma}\), and its supremum is over the published
uncertainty-class intersection represented by
\(\mathcal P_{\mathrm{ACE},n}^{\mathrm{JMS}}\).  Combining (3) and (4)
also gives the restricted inequality on the full spectral class.  This
proves every assertion of the proposition.